\section{Why the channels differ}
\label{sec:why-channels-differ}

A belief in context is a sentence the model can quote, reconsider, and abandon. A
compiled belief conditions the forward pass without ever appearing in it. Three
measurements are consistent with that account, though none of them proves it outright.

The clearest evidence comes from splitting the modulation by what it does to two
different projections. We compiled the same 24 belief trees on both sides of each
topic and measured the functional similarity of the resulting query- and value-deltas
across belief sets, on a fixed probe set. The query delta is nearly identical no matter
which belief produced it, mean off-diagonal cosine 0.996 across belief sets (pro side,
\(n=24\)), a content-blind piece of machinery that looks the same whether the compiled
belief is about space colonization or encryption backdoors. The value delta tracks what
the belief actually says: mean off-diagonal cosine 0.53, with a wide spread
(0.18--0.92).\expref{app:exp-e1} This split replicated on a held-out con-side confirmation set,
pre-registered before the run: dQ 0.996, dV 0.53 again, both legs of the prediction
holding independently of which side of a topic the model
argued.\expref{app:exp-e1b} Query modulation
supplies a content-blind capacity to hold; value modulation supplies the content it
holds. The paper's dissociation between installing and holding shows up a second time,
inside the mechanism, between two projections of the same update.

A linear probe trained to decode the active belief from hidden states supports the same
reading from a different angle: it reaches 100\% accuracy when modulation is active and
22\%, below chance, when modulation is zeroed, so the hidden states carry no
recoverable information about which belief is compiled once the modulation that
carries it is removed, consistent with attention-pattern divergence between modulated
and unmodulated passes being small (Jensen--Shannon divergence 0.049). Modulation also
measurably suppresses think-mode generation, the channel through which the model would
otherwise narrate, and thereby expose to an interlocutor, its own reasoning about
whether to concede, an effect that holds regardless of rank or
input-dependence\iflongversion{} (Section~\ref{sec:isolation})\fi.

The two channels are dissociable, and the compiled one affects processing without
entering the reasoning trace the model produces; that is the claim the evidence above
supports, no further. Text in context pays two costs, not one. It is available to be
quoted and argued against, and a length-matched control shows it is also diluted simply
by being long: belief text replaced word for word with topic-blind filler at the same
token count holds at 55.6\%, a point estimate below the real belief's own 58\%
(Section~\ref{sec:dissociation-prior}; the confidence interval at this sample size does
not cleanly separate the two accounts). Arguability and length are both intrinsic to
what it costs to store a belief as text a model has to keep re-reading against a
growing transcript, not competing explanations for why context fails. Compiled
modulation pays neither cost: it is not quotable, and it does not occupy positions that
dilute.
